% What Scale Buys in Tool-Use Agents, and How to Buy It Cheaply.
% Compile: pdflatex what_scale_buys ; bibtex what_scale_buys ; pdflatex x2. Compiles on Overleaf.
% To target a venue, replace the geometry preamble with the venue style file.
\documentclass[10pt]{article}
\usepackage[letterpaper,margin=1.7cm]{geometry}
\usepackage{times}
\usepackage[T1]{fontenc}
\usepackage[utf8]{inputenc}
\usepackage{textcomp}
\usepackage{microtype}
\usepackage{booktabs}
\usepackage{array}
\usepackage{amsmath,amssymb}
\usepackage[numbers,sort&compress]{natbib}
\usepackage[hidelinks]{hyperref}

\title{\textbf{What Scale Buys in Tool-Use Agents, and How to Buy It Cheaply:\\
A Capability\,$\times$\,Scale\,$\times$\,Lever\,$\times$\,Cost Map Where Compliance Is Scale-Invariant}}
\author{Anonymous submission}
\date{}

\begin{document}
\maketitle

\begin{abstract}
Tool-use agents map a natural-language request to a sequence of tool calls under a domain policy, and are usually improved by scaling the base model. We ask instead, function by function and across scale, what scaling actually buys: which sub-capabilities improve with size, which remain limited at every scale, which are better supplied by deterministic scaffolding or by per-domain authored relations (which we call A2), and how these should be assembled to minimize total deployment cost. Treating four benchmarks ($\tau^2$-bench, SOPBench, TaskBench, and Synth) as measurement instruments, we decompose tool-use failure into \emph{operand execution}, \emph{orchestration load}, and \emph{compliance}, and chart each across model scale (7B--32B measured; 72B and above projected) and across levers (base model, capability scaffold, guarantee scaffold, A2, and learning). We report three findings. First, operand execution saturates early: given an explicit specification, a 32B base model selects the correct action on 88 of 88 cases, so this is not a learnable gap at that scale. Second, the load dimensions that predict failure retire with scale in a definite order---interference and state-multiplicity bind at 7B and 14B but disappear by 32B, while context length and conditional depth persist---and a controlled probe overturns an interference signal that appears only in observational correlations. Third, compliance is scale-invariant: task success rises with scale (0.19 at 7B to 0.55 at 32B), yet the per-write-opportunity rate of policy-semantic violations does not fall ($\approx$0.07--0.10 across all measured scales, with overlapping 95\% confidence intervals), even as authorization-type violations collapse; a deterministic policy gate eliminates the residual at every scale. We interpret deterministic scaffolding as a load-reduction operator and fold the map into a total-lifecycle cost objective. On $\tau^2$-retail, a 32B int8 base with a deterministic gate runs on-premises at roughly an order of magnitude lower cost per request than a frontier API at near-equal compliance, and a small-on-premises-plus-escalation fleet is projected to reach equal-or-better compliance at lower cost. We conclude that scale buys capability but not guarantee: even a frontier model requires a deterministic guarantee layer, so the cost-optimal design is not the smallest model but the total-cost knee.
\end{abstract}

\section{Introduction}
A tool-use agent maps a natural-language request to a sequence of tool calls that satisfy a domain policy. The default recipe for improving such an agent is a larger base model. But ``larger'' conflates distinct capabilities, and it is expensive on two permanent axes: inference cost (a model twice the size costs more than twice as much per request and is slower, for the lifetime of the deployment) and, for on-premises or sovereign deployment, hardware capital expenditure, since the operator must buy the memory, power, and cooling. The practical question is therefore not \emph{how large a model} but \emph{which capabilities must scale, which can be delegated to cheap deterministic structure, which no model can be trusted to perform, and what mixture minimizes total cost}.

\paragraph{Two deployment realities motivate this study.} First, regulated and closed deployments force small on-premises models and make compliance a first-class concern. In regulated or sovereign settings---finance, healthcare, government---customer data cannot leave the premises, so a frontier API is unavailable and the operator runs, and buys the hardware for, a local model. At the same time, regulation increasingly demands traceability, validation, and effective human oversight (for example EU AI Act Articles 12 and 14, and US banking guidance SR~26-2)---but not determinism as such; no major framework textually mandates it, the medical-device provision EU MDR Annex~I~\S17.1 being the lone exception. A deterministic policy gate satisfies these obligations by construction, and under SR~26-2 deterministic rule-based processes fall outside the model-risk perimeter entirely, whereas a stochastic agent can demonstrate compliance only statistically, per sample, with its effective policy unverifiable in advance. Our central finding---that compliance is scale-invariant (\S\ref{sec:compliance})---implies that such a gate is needed at every scale, frontier models included.

Second, per-domain authoring is the recurring cost, so a fixed scaffold combined with cheap-to-author relations is the path to many domains. Domain-specific learning is slow and expensive and must be redone for each field; if the per-domain artifact (the A2) can be authored cheaply and the scaffold and learned skill are held fixed, then moving to a new domain costs an A2-swap rather than a retrain. This paper measures single-domain compliance and cost; the multi-domain payoff is the generalization axis of the cost framework (\S\ref{sec:knee}), which we validate only partially here (a SOPBench cross-domain check) and otherwise treat as a design premise rather than a measured multi-benchmark result. Taken together, a deployer must use a small on-premises model for cost and sovereignty, needs a form of compliance that scale does not provide, and wants to amortize effort across domains; the map developed here speaks to all three.

We answer with a capability $\times$ scale $\times$ lever $\times$ cost decomposition. For each functional sub-capability we ask, across scale, whether it (a)~grows with scale, (b)~is a genuine model limit, (c)~is covered by deterministic \emph{scaffold}, or (d)~is covered by per-domain authored relations (\emph{A2}); we then assemble the per-capability levers under a total-lifecycle cost objective and locate the cost knee. The benchmarks are instruments for filling these cells, not the object of study.

\paragraph{Objective.}
We seek to minimize, simultaneously, the cost of cheap GPU tier and count, the memory footprint (favoring int8 or int4 quantization), the maintenance burden per change, the number of human experts required, and the cost of generalizing to each new field (ideally an A2-swap alone), together with hardware capital and inference operating expenditure. These terms are coupled and point in the same direction: cheap, low-memory hardware favors a small, quantized model, while low maintenance and onboarding favor maximizing the fixed domain-general component (deterministic scaffold plus transferable, forgetting-free learning) and minimizing the per-domain A2. The A2 is the common adversary, since it drives maintenance, experts, and per-field cost, and is worst on-premises where a frontier A2-generator is unavailable; a fixed scaffold and transferable skill are the common ally. The optimum is the total-cost knee, not the smallest model.

\paragraph{Contributions.}
(1)~A controlled, frontier-API-free methodology that isolates operand execution, orchestration load, and compliance across scale (\S\ref{sec:method}).
(2)~Evidence that operand execution is scale-saturated and not a learnable gap at 32B (\S\ref{sec:operand}); the apparent gap is criterion interpretation and join/disambiguation. Faithful-formalization fine-tuning is therefore unwarranted at 32B, though a residual may exist at 7B (\S\ref{sec:future}).
(3)~A load decomposition with a scale response (\S\ref{sec:load}); controlled generation overturns a correlation visible only in observational data.
(4)~The finding that compliance is scale-invariant (\S\ref{sec:compliance})---the empirical backbone of a two-kind scaffold (a capability scaffold that recedes with scale, and a guarantee scaffold that does not).
(5)~A total-cost-optimal lever allocation with a deployment payoff (\S\ref{sec:cost}--\ref{sec:knee}): decidability-first routing, an on-premises cost-per-request advantage of roughly an order of magnitude (the exact multiple reported as an estimate pending measured token accounting), and a projected fleet design, framed by the cost knee.
(6)~A proposed unifying operator---deterministic scaffold as load reduction (\S\ref{sec:framework})---for which we establish the components but defer the full predictive test (\S\ref{sec:future}), and which we therefore present as a framework hypothesis rather than a verified law.
Our claim is deliberately \emph{not} ``small beats large,'' a result already established by reinforcement-learned routing systems (\S\ref{sec:related}); it is \emph{which} parts of the gap are scale-invariant and must be offloaded, and a cost-optimal map indicating which lever a deployer should pull for each capability.

\section{Related Work}\label{sec:related}
The central lineage is conceptual---why a language model is good at language yet must offload deep procedures and cannot self-guarantee compliance---rather than the recent ``small-beats-large'' routers.

\paragraph{Language is shallow-parallel; deep procedures must be offloaded.} A bounded-depth Transformer forward pass is asymptotically uniform TC$^0$ \citep{ax2207_00729}; chain-of-thought escapes this only by serializing depth \citep{ax2305_15408,ax2402_12875}. Natural language embeds procedures of definite computational class---van Benthem's semantic automata (the quantifier ``most'' is pushdown), which are cognitively real; Relevance Theory's procedural/conceptual distinction; Goodman's notationality, which concerns cost rather than computability (Zhang \& Norman). \emph{Classifying} which deep procedure a request embeds is shallow and learnable, whereas \emph{executing} it is serial-deep and best offloaded---the basis of our capability/lever routing.

\paragraph{Cognitive architectures, non-introspective decidability, scale-invariant epistemic limits.} We instantiate the CoALA umbrella \citep{ax2309_02427}, making concrete what it leaves general: a relational-algebra deterministic core, empty-relation abstention, and zero-retrain A2-swap. SOAR \citep{ax2505_07087} surfaces indecision architecturally---an \texttt{impasse} arising from an absent preference rather than from introspection---the precedent for our pattern of mapping uncertainty to an observable empty or ambiguous relation and then asking. Deterministic-first with language-model fallback is established \citep{ax2403_00810,ax2510_09355}. We add a formalization-faithfulness axis, a \emph{learned} empty-to-ask decision, and a small weight-learned domain-general skill with A2-swap. Hallucination is provably scale-invariant \citep{ax2401_11817,ax2509_04664}; our \S\ref{sec:compliance} is the policy-enforcement analogue---no model self-guarantees compliance, just as none self-guarantees truthfulness, and both demand an external deterministic mechanism.

\paragraph{Language model proposes, deterministic component executes.} LLM-Modulo \citep{ax2402_01817}, PAL \citep{ax2211_10435}, natural-language-to-PDDL formalizers, and reinforcement learning from verifiable rewards with deterministic checkers (T\"ulu~3, \citealp{ax2411_15124}) all follow this pattern. The boundary is honest: offload does not always win, owing to formalization overhead on simple cases and to specification leakage, and we treat that boundary as itself measurable.

\paragraph{Guarantees and shielding.} Safe-reinforcement-learning shielding \citep{ax1807_06096} and predictive safety filters, and their agent-level instances ShieldAgent \citep{ax2503_22738}, AgentSpec \citep{ax2503_18666}, and Formal-LLM \citep{ax2402_00798}, veto actions using hand-written or probabilistic specifications. Our gate differs in being deterministic, in \emph{driving} (repairing) preconditions rather than only vetoing, and in pairing with a learned proposer and zero-retrain transfer.

\paragraph{Learned routing and small-model orchestration (the closest rivals).} Routing is a learned, transferable skill (A2FM \citep{ax2510_12838}, xRouter \citep{ax2510_08439}, When2Call \citep{ax2504_18851}); tool-necessity is even linearly decodable from hidden states (AUROC 0.89--0.96, \citealp{ax2605_09252}). ToolOrchestra \citep{ax2511_21689} reinforcement-learns an 8B router that delegates to a frontier model, reaching 80.2\%@10.3\textcent\ versus GPT-5's 77.7\%@31.3\textcent\ on $\tau^2$---the strongest cost result to date, but one that relies on frontier delegation and is therefore disqualified in the on-premises, no-egress regime; it optimizes cost given black-box components and offers no deterministic offload, by-construction compliance, decidable-versus-learned accounting, or zero-retrain transfer. TRUST \citep{ax2606_06976} aligns tool-calling decisions to model uncertainty by reinforcement learning but remains black-box. ReDAct \citep{ax2604_07036} and unified routing-cascading \citep{ax2410_10347} perform calibrated deferral; our fleet is a deterministic-router instance of the same idea. These works establish that a cost-favorable small-model result exists, which is precisely why our contribution is the analysis and method---what is scale-invariant, which lever fits each capability, and where the cost knee lies---rather than the existence result itself.

\paragraph{Cost, regulation, and industry.} Against a pure-scale reading of the Bitter Lesson, Brooks' ``A Better Lesson'' (2019) and Chollet \citep{ax1911_01547} argue that total cost and priors must be counted; our knee operationalizes this. The buyer value of determinism is verifiability and model-risk exemption (SR~26-2; EU AI Act 2024/1689; EU MDR Annex~I~\S17.1 on repeatability). Palantir Foundry, whose Ontology resembles scaffold plus A2 and whose AIP language-model layer resembles a small model with a transferable skill, is the industrial baseline; we differentiate on cost---hardware, experts, configuration-only field-crossing, and zero egress---rather than on functionality.

\paragraph{Load and benchmarks.} The novelty here is only the combination of load decomposition with scaffold-as-load-reduction for tool use. We use benchmarks as instruments: $\tau^2$-bench \citep{ax2506_07982} (the $\tau$-bench predecessor, \citealp{ax2406_12045}), SOPBench \citep{ax2503_08669}, TaskBench, and Synth/CFB \citep{ax2501_10132}; a precedent for measuring a capability floor in this way is the Sufficient Context study (Joren et al., ICLR 2025).

\paragraph{The open intersection.} Prior work supplies each piece---procedural semantics and automata, Transformer expressivity, SOAR and CoALA, hallucination inevitability, formalize-and-offload, learned routing, and shielding. The unoccupied intersection is the empirical capability $\times$ scale $\times$ lever $\times$ cost map: which capabilities scale retires, which are genuine scale-invariant limits (compliance, measured here), which lever covers each, and the total-cost-optimal assembly. Recent routers optimize cost over black boxes; we open the box.

\section{The Capability$\times$Scale$\times$Lever$\times$Cost Framework}\label{sec:framework}
For each function $f$ we populate a row indexed by scale $N$, lever $\ell \in \{\text{base, capability scaffold, guarantee scaffold, A2, learning}\}$, and assembled cost. The framework rests on three structural claims.
\begin{itemize}\itemsep1pt
\item \textbf{Operand vs.\ orchestration vs.\ compliance} fail for different reasons and respond to different levers: operand execution is atomic execution given a criterion; orchestration is holding interdependent steps coherent under load; compliance is obeying policy with a guarantee.
\item \textbf{Decidability-first lever routing} (cost-ascending short-circuit): decidable from schema, policy, and state $\to$ scaffold; else a domain fact $\to$ A2; else promptable at the available size $\to$ prompt; else scale-emergent and affordable $\to$ scale or replicate cheaply; else an irreducible natural-language-to-formal step $\to$ a minimal transferable adapter (last resort). This inverts the ``scale it'' reflex.
\item \textbf{Two kinds of scaffold.} A \emph{capability scaffold} (presenting records, computing a value, resolving a reference) substitutes for a missing skill and recedes as scale grows (F3, raw pass). A \emph{guarantee scaffold} (policy gates) supplies a property no probabilistic model provides and is scale-invariant (F4, compliant pass). Section~\ref{sec:compliance} measures this split---why even a frontier system retains a deterministic trust layer.
\end{itemize}
\textbf{Load as a vector.} Orchestration load is not scalar: $L_{\mathrm{len}}$ (context to retain), $L_{\mathrm{state}}$ (interdependent state), $L_{\mathrm{branch}}$ (conditional depth), $L_{\mathrm{interf}}$ (confusable entities), $L_{\mathrm{contra}}$ (mid-stream revision), with load orthogonal to operand difficulty. \textbf{Scaffold as a load-reduction operator:} each scaffold removes a dimension (present/autofetch lower $L_{\mathrm{len}},L_{\mathrm{interf}}$; compute lowers $L_{\mathrm{state}}$; gates lower $L_{\mathrm{branch}}$; a controller lowers $L_{\mathrm{state}},L_{\mathrm{len}}$), giving $L_{\mathrm{eff}} = L - \Delta L_{\mathrm{scaffold}}$ and the prediction $\mathrm{pass}(f,N) \Leftrightarrow L_{\mathrm{eff}}(f) < L^*(N)$, where $L^*(N)$ is the load tolerance at scale $N$. To avoid tautology, $\Delta L$ should be estimated independently (from the mechanical feature reduction the scaffold performs) and shown to predict the shift in failure onset. We establish the components here and defer the full predictive test above 32B to \S\ref{sec:future}, presenting the operator as a hypothesis rather than a validated law.

\section{Method}\label{sec:method}
\textbf{Benchmarks as instruments.} $\tau^2$-retail is primary; SOPBench, TaskBench, and Synth are the training substrate for the lever allocation, while $\tau^2$ is held out as a transfer target and is never trained on. Scale axis: Qwen2.5 at 7B, 14B, and 32B (GPTQ int8), all measured; 72B and 235B projected (\S\ref{sec:future}).
\textbf{Cost discipline.} The only paid component of the evaluation is a GPT-4.1 user simulator, which we avoid during discovery; capability is measured with isolated single-shot probes, existing multi-scale runs, and a local user simulator (faithful scripted turns over a local base model) for full-flow checks. Paid full runs serve only as a final confirmation of a conclusion already reached without them, at minimal scope.
\textbf{Conditions.} A \emph{floor} condition (no scaffold) versus capability- and guarantee-scaffold conditions (\emph{g15}, \emph{present}, \emph{nested}, \emph{assembled}); a \emph{given-spec} versus \emph{goal} manipulation separates operand execution from criterion interpretation.
\textbf{Metrics.} A robust pass rate (pass over all repeated trials per item; the user-simulator noise is $\approx0.11$, so single-trial figures are not reported as findings); a raw task-pass rate F3; a compliant-pass rate F4 (F3 less policy violations); per-dimension partial correlation of failure with each load feature (controlling for operand difficulty); controlled accuracy-versus-load onset curves; and a cost row per cell (latency, tool round-trips, memory and GPU tier, cost per request).

\section{Results}
\subsection{Operand execution is scale-saturated, not a learnable gap at 32B}\label{sec:operand}
Varying only the user-simulator input: given an explicit specification (\textbf{given-spec}) the 32B base is \textbf{88/88 (100\%)}; given only the goal (\textbf{goal}) it is 62/88 (70\%), the gap being criterion interpretation (argmax/argmin computation, multi-attribute reasoning, conversational fidelity) rather than execution. Genuine join and disambiguation is 7/13 (54\%) and is addressable by presenting the candidate records. Operand execution is therefore not a 32B capability gap, and faithful-formalization fine-tuning is unwarranted at that scale. (Every earlier ``operand failure'' in our forensics dissolved into a measurement artifact---a computation bug, a premature case split, or a probe that supplied the identifier---until the controlled probe read 100\%. The cost corollary: at 7B the base may be genuinely weak even given a specification, so the learning lever can re-acquire a role at small scale; see \S\ref{sec:future}.) A complementary limit: large-candidate selection does not scale in-head---comparative selection at 50 candidates scores 0.02 and ranking at 50 scores $\approx0.30$ even including a 235B model, whereas a deterministic engine scores 1.00. (The 235B figure appears only in this large-candidate probe; the scale-axis compliance and load runs span 7B--32B, with 72B and 235B projected.)

\subsection{Capability-type load retires with scale, in a definite order}\label{sec:load}
Partial correlation of \emph{floor} failure with each load feature (controlling for operand difficulty):
\begin{center}\small\begin{tabular}{lccc}
\toprule dimension & 7B & 14B & 32B \\ \midrule
$L_{\mathrm{len}}$ (length) & $+0.20$ & $+0.33$ & $\mathbf{+0.37}$ \\
$L_{\mathrm{state}}$ & $\mathbf{+0.24}$ & $\mathbf{+0.21}$ & $+0.06$ \\
$L_{\mathrm{branch}}$ (conditional) & $+0.20$ & $+0.13$ & $\mathbf{+0.26}$ \\
$L_{\mathrm{interf}}$ & $\mathbf{+0.30}$ & $\mathbf{+0.20}$ & $+0.08$ \\
$L_{\mathrm{contra}}$ & $+0.06$ & $+0.13$ & $+0.04$ \\ \bottomrule
\end{tabular}\end{center}
The binding set is scale-dependent: four dimensions at 7B/14B, only length and conditional depth surviving by 32B (interference and state retire first). $L_{\mathrm{contra}}$ is too sparse in $\tau^2$ to fit reliably. A controlled interference probe (fix the operand; vary only the number of confusable variants) does \emph{not} reproduce the strong 7B signal:
\begin{center}\small\begin{tabular}{lccccc}
\toprule $N$ & 1 & 2 & 4 & 8 & 16 \\ \midrule
32B & 1.00 & 1.00 & 1.00 & 0.97 & 0.87 \\
7B  & 1.00 & 1.00 & 1.00 & 0.90 & 0.80 \\ \bottomrule
\end{tabular}\end{center}
the gap at $N=16$ is only $0.80$ versus $0.87$, so the observational interference correlation was largely a size confound. (Single-shot probes under-test the inherently multi-turn dimensions $L_{\mathrm{state}}$ and $L_{\mathrm{contra}}$; controlled multi-turn probes are left to future work.)

\paragraph{The residual is orchestration, not operand.} Under faithful local turns, 2/10 of the robust fail-all set are simulator-noise flips; the dominant residual is orchestration under load (precondition-violating batching; multi-order tracking that drops an order and fabricates an identifier; conditional sequencing; context overflow). An isolation plan probe (eliciting the abstract plan with all reads supplied, grading only structure) yields correct structure on 6/10---roughly half execution load (closable by separating planning from execution deterministically) and roughly half an in-isolation planning miss (the candidate for a learned path-selection lever, left to future work). The assembled-stack ceiling is a robust pass rate of $0.402$.

\subsection{Compliance is scale-invariant --- the central finding}\label{sec:compliance}
Separating F3 (task pass) from F4 (compliant pass) on \emph{floor} (no gate):
\begin{center}\small\begin{tabular}{lccccc}
\toprule scale & F3 & F4 & gap & violations (g1 auth / g2 policy) & total viol.\% \\ \midrule
7B  & 0.189 & 0.130 & 0.059 & 57 / 31 & 33.6\% \\
14B & 0.468 & 0.404 & 0.064 & 38 / 38 & 26.3\% \\
32B & 0.547 & 0.491 & 0.056 & 6 / \textbf{41} & 14.3\% \\ \bottomrule
\end{tabular}\end{center}
The absolute F3--F4 gap is roughly flat (0.059/0.064/0.056), and policy-semantic violations (g2) do not fall ($31\to41$) even as authorization-type violations (g1) collapse ($57\to6$). With the gate (g15): F3${=}$F4, gap ${=}0.000$, violations ${=}0.0\%$ at every scale. Scale solves capability-type compliance (g1) but not policy-semantic compliance (g2); the gate zeros the residual at every scale, frontier included---the policy-enforcement analogue of hallucination inevitability and the backbone of the two-kind scaffold.

\paragraph{Resolving the completion-rate confound.} The absolute g2 count ($31\to41$) and total violation percentage ($33.6\to14.3$) are confounded by completion: a higher-pass model completes more trajectories and attempts more writes ($526\to683\to680$ across 7B/14B/32B), so raw counts conflate ability with exposure. Normalizing to the per-write-opportunity rate removes this confound: the g2 rate per write opportunity is $0.103$ (95\% Wilson interval $[0.080,0.132]$) at 7B, $0.070$ ($[0.053,0.092]$) at 14B, and $0.075$ ($[0.058,0.097]$) at 32B---flat across scale, with overlapping confidence intervals and no significant downward trend (if anything the 7B rate is marginally higher). Scale therefore does not reduce the confirm-before-write violation rate ($\sim$1 in 10--14 writes at every scale), while the gate zeros it everywhere; the strong scale-invariant form holds, not merely the differential (authorization-type violations still collapse while policy-semantic violations persist). This range is 7B--32B; the projected 72B point (\S\ref{sec:future}) would extend it. This result currently rests on one violation class (confirm-before-write) on one benchmark, a scope limit we return to in \S\ref{sec:limits}.

\subsection{Cheap replication: a capability scaffold substitutes for scale}
Raw task pass rises monotonically (single-trial pass $=$ 7B 0.24 / 14B 0.52 / 32B 0.60, $n{=}342$). (This single-trial pass rate comes from a separate scale-decomposition run and is a different estimator from the \S\ref{sec:compliance} floor-compliance F3 of $0.189$ to $0.547$, which is robust over three trials; both measure raw task pass, differing by estimator and run, not by contradiction.) Each piece is obtainable more cheaply than scale: a deterministic fetch-first engine (deny ungrounded provenance, call the producer, inject the real value) moves grounding errors $33\to9$ and roughly doubles pass ($0.14\to0.264$) with no learning. A lever audit: authorization, confirmation, ownership, and precondition gates are genuine levers, whereas an eligibility-\emph{steering} gate has essentially no causal effect and naive retry is counterproductive.

\subsection{The map, filled}
\begin{center}\small\begin{tabular}{p{3.1cm}p{2.8cm}p{1.4cm}p{4.4cm}}
\toprule function & scale response & limit? & efficient lever \\ \midrule
operand execution & saturates by 32B & no (at 32B) & base; learning may re-acquire a role at 7B \\
interference / state load & retire by 32B & no & capability scaffold (present), or scale \\
length / conditional load & persist at 32B & partial & capability scaffold (partial); conditional $\to$ controller or learn \\
large-candidate selection & flat at all scales & yes (in-head) & deterministic computation \\
policy compliance (g2) & \textbf{scale-invariant (per-write rate flat)} & \textbf{yes (guarantee)} & \textbf{guarantee scaffold (gate)} \\
domain facts & --- & --- & A2 / retrieval \\ \bottomrule
\end{tabular}\end{center}

\subsection{Cost: the deployment payoff}\label{sec:cost}
On $\tau^2$-retail, using existing runs, a 32B int8 base with a deterministic gate runs on-premises at $\approx\$0.0019$ per request versus a frontier API at $\approx\$0.044$ per request, a ratio of about 23$\times$ (and 16--40$\times$ over a GPU priced at \$0.2--0.5 per hour); it is $\sim$3000$\times$ cheaper than a human agent at $\sim\$6$ per contact, and is auditable with zero data egress. The trade-offs are honest: latency is $\sim$6$\times$ higher (178s versus 30.5s), and pure on-premises compliance (0.573) is below the frontier value (0.82), so the claim is near-equal compliance at $\approx$23$\times$ lower cost, not equal compliance. (The 0.573 figure is gated 32B compliance from the deployment run at a single trial, distinct from \S\ref{sec:compliance}'s 0.547 over three trials---the same quantity on a different trial count and run. Cost per request is estimated from a token count of approximately characters-over-four with an assumed GPU utilization; the ratio is robust to an order of magnitude, but the exact multiple awaits measured token accounting, \S\ref{sec:future}.)

A heterogeneous fleet is projected to improve on this: routing easy requests to the 32B on-premises model and hard requests to a frontier model yields a blended compliance of 0.860 (above the pure-frontier 0.816) at \$0.021 per request, $\sim$2.1$\times$ cheaper, on clean data. We report this as a projection, not a measurement: an arithmetic upper bound that assumes a cheap decidable router. Such a router is feasible---it can reuse the decidable gate and uncertainty mechanism of \S\ref{sec:compliance} with no new learned component---but is not yet implemented.

\section{Cost-Optimal Lever Allocation and the Knee}\label{sec:knee}
The map yields a deployment rule that inverts the ``scale it'' reflex: reduce the load and supply the guarantee with structure, rather than training the model to tolerate either. A capability scaffold removes the load dimensions scale would otherwise have to buy, and recedes when scale is available; a guarantee scaffold supplies the property scale never buys, and is retained at every scale. Per-capability allocation is decidability-first (\S\ref{sec:framework}); the A2 is the common adversary (minimized) and a fixed scaffold plus transferable skill the common ally. The optimum is the cost knee---the smallest total-cost size---and, more generally, a heterogeneous on-premises fleet whose mixture is set by the per-capability crossover between learning and an engine. Where a load is irreducible by scaffold (conditional depth and path search are our candidates), the lever shifts to a learned, transferable heuristic rather than a larger model. The positioning relative to Palantir Foundry is on cost (hardware, experts, field-crossing, egress), not on functionality.

\section{Future Work}\label{sec:future}
Several cells of the map are measured here and several remain projected; the natural next steps, in order of value, are as follows. The most consequential is a compliance measurement at 72B and above, which would extend the F3/F4 scale-invariance result beyond 32B---the single largest credibility gain for the central claim. Second, the compliance result should be replicated across additional policy-violation classes and additional benchmarks (SOPBench in particular), since it currently rests on one violation class on one benchmark. Third, controlled multi-turn probes for the state and contradiction load dimensions, and controlled onset curves for length and conditional depth, would complete the load decomposition. Fourth, an independent estimate of $\Delta L$ and a test of the load-reduction prediction at two or more scales would move the scaffold-as-load-reduction operator from hypothesis to validated relation. Fifth, a cross-field A2-swap (for example to airline or banking domains), with the fixed component held invariant and no retraining, would upgrade transfer from a design premise to a partially measured result. Finally, enabling measured token accounting would turn the cost-per-request figures from estimates into measurements and allow a sweep of size against total cost to locate the knee and the fleet mixture.

\section{Discussion}
The result reframes the deployment question from model size to lever allocation under cost. Two asymmetries do the work: capability is scale-retired \emph{in a definite order} (a recede-able capability scaffold buys a small model the dimensions scale would otherwise provide), and compliance is \emph{scale-invariant} (a guarantee scaffold is retained even at the frontier). Training a model to tolerate load is training it to do the deterministic controller's job; the design instead routes that work to determinism by default and to learning only for the residual---a residual that is itself scale-dependent (empty for operand execution at 32B, non-empty at 7B), which is why the cost knee, rather than a fixed model size, is the right object.

\section{Limitations}\label{sec:limits}
\textbf{Scale range:} 7B--32B measured; 72B/235B/frontier projected (\S\ref{sec:future}). The scale-invariance claim is a flat trend supported by priors and would be strengthened by the 72B measurement. \textbf{Compliance metric:} the strong scale-invariant form is supported by the per-write-opportunity g2 rate (flat across 7B--32B, overlapping CIs; \S\ref{sec:compliance}), not only by the differential, but rests on one violation class (confirm-before-write) on one benchmark; cross-class and cross-benchmark replication and the 72B point remain to be done. \textbf{Observational versus controlled:} \S\ref{sec:load} correlations are size-confounded; only the controlled interference probe is causal, and controlled length/conditional curves and multi-turn state/contradiction probes remain to be run. \textbf{Cost figures:} cost per request is token-estimated and robust to an order of magnitude but not yet measured; the fleet result assumes a decidable router not yet built; the human and Palantir comparisons are external baselines. \textbf{Transfer:} zero-retrain A2-swap is a design premise; cross-field validation is partial, and the mechanism is itself a variant of schema-guided dialogue state tracking and tool-use systems---we claim the scale-invariance map, not the transfer mechanism. \textbf{Single benchmark for compliance:} the F3/F4 analysis is on $\tau^2$-retail.

\section{Conclusion}
Scaling a tool-use agent buys capability---it retires load dimensions in a definite order and lifts task pass---but it does \textbf{not} buy compliance: policy-semantic violations are flat across 7B$\to$32B even as pass climbs, and a deterministic gate removes them at every scale. The right object of optimization is not the model size but the capability$\times$scale$\times$lever$\times$cost map: a small base, a capability scaffold that recedes as scale grows, a scale-invariant guarantee scaffold, authored A2, and learning reserved for the scale-dependent residual, assembled at the total-cost knee. This is why even frontier systems should use deterministic scaffolding and authored relations, and it motivates the companion work that generates the A2 and learns the path-selection residual.

\small
\bibliographystyle{plainnat}
\bibliography{references}
\end{document}
